\section{Results and Discussion} \label{sec:results}

% \input{fig/dev-sequencescreen}
% \input{fig/dev-sequencescreen-neutral}
% \input{fig/dev-sequencescreen-pessimal}
% \input{fig/dev-sequencescreen-beneficialinhost}

In this section, we report results of trials testing the efficacy of proposed extended methods for screening for multilevel selection effects against ground-truth data generated \textit{in silico}.
We discuss the fitness effect screen component performance first, followed by discussion of the mutation frequency screen component.

\subsection{Example Fitness Effect Screen}

\begin{figure}

\centering
\includegraphics[width=0.85\linewidth]{binder/binder-2025-05-17-vanilla-compscreen.ipynb/binder/teeplots/col=trt-name+hue=metric+kind=boxen+row=comparator-set+viz=catplot+x=metric+y=percentile+ext=}

\caption{%
\textbf{Sample fitness effect screen result, baseline conditions.}
\footnotesize
Boxen plots represent distributions of percentile-normalized clade fitness metric values for focal site MDCs, with strip plot dots corresponding to individual MDCs.
Vertical bars indicate bootstrap 95\% confidence interval; annotated $p$ values correspond to bootstrap comparison of sample mean against null expectation of mean equal to 50th percentile.
Bottom row shows results for MDC comparator sets partitioned by simulation month; top row comparator sets for percentile ranking included all observed MDCs.
}
\label{fig:results-vanilla-example}

\end{figure}


As a first step in evaluating our comparator-based approach for assessing MDC fitness outcomes, we compared the results of sample assays for deleterious between-host transmission effects across Gneu/Sneu, Gneu/Sben, and Gdel/Sben treatments.
For illustrative purposes, we show as examples the strongest-surveyed 2$\times$ magnitude for Gdel and Sben effects.

Figure \ref{fig:results-vanilla-example} compares the distribution percentile-normalized clade fitness statistics for focal MDCs across the neutral, within-host advantage, and true multilevel selection tradeoff treatments.
For multilevel selection Gdel/Sben conditions (Figure \ref{fig:results-vanilla-example:TODO}), bootstrap testing indicated the mean percentile ranking of MDC clade sizes and durations to both be significantly below the 50th percentile ($\mu=TODO$, $p<TODO$; $\mu=TODO$, $p<TODO$).
TODO is this with within-month normalization?
Promisingly, significant evidence of MDC fitness effects were not detected under either the Gneu/Sneu or Gneu/Sben control conditions.
For Gneu/Sneu conditions, note that too few mutations are observed for the bootstrap test to be considered reliable.

\subsection{Fitness Effect Screen}

\begin{figure*}
\begin{minipage}{0.65\linewidth}
\includegraphics[width=\linewidth, trim={0 0 0.95cm 0}, clip]{binder/binder-2025-05-19-compscreen-vanilla-summary.ipynb/binder/teeplots/2025-05-19-compscreen-vanilla-summary/col=trt-name+comparator=month+hue=sig+kind=swarm+row=y+viz=catplot+x=smallsamp+y=p-boot+ext=}
\subcaption{fitness effect screen covasim population size 200,000}
\label{fig:vanilla-summary-results:fit-effect-small}
\end{minipage}%
\begin{minipage}{0.35\linewidth}
\includegraphics[width=0.497\linewidth, trim={0.6cm 0 0.91cm 0}, clip]{binder/binder-2025-05-19-compscreen-vanilla-big-summary.ipynb/binder/teeplots/2025-05-19-compscreen-vanilla-big-summary/col=trt-name+comparator=month+hue=sig+kind=swarm+row=y+viz=catplot+x=smallsamp+y=p-boot+ext=}
\includegraphics[width=0.483\linewidth, trim={1.65cm 0 0 0}, clip]{binder/binder-2025-05-19-compscreen-vanilla-big-1.3x-summary.ipynb/binder/teeplots/2025-05-19-compscreen-vanilla-big-1.3x-summary/col=trt-name+comparator=month+hue=sig+kind=swarm+row=y+viz=catplot+x=smallsamp+y=p-boot+ext=}
\subcaption{fitness effect screen covasim population size 1.2 million}
\label{fig:vanilla-summary-results:fit-effect-big}
\end{minipage}

\caption{%
\textbf{Fitness effect screen results.}
\footnotesize
TODO
}
\label{fig:vanilla-summary-results}
\end{figure*}


Given the promising results from example screens shown in Figure \ref{fig:results-vanilla-example}, we next sought to assess detection outcomes across replicate experiments and smaller effect $1.1\times$ and $1.3\times$ effect sizes.
Figure \ref{fig:vanilla-summary-results} summarizes bootstrap tests for depression of mean normalized clade size and duration associated with focal MDCs.
Experiments were conducted under baseline ``vanilla'' configuration.

In a first set of assays, we tested a naive approach without comparator MDC partitioning.
That is, every MDC was included in the comparator set for all MDCs.
However, we found this approach to produce a high false-positive rate under Gneu/Sneu and Gneu/Sben conditions (TODO/TODO trials at $\alpha=0.05$).
We found that binning comparator sets by simulation months was effective in ameliorating this defect in test specificity, and so adopted the approach of partitioning comparator sets by month for all subsequently-discussed results.

Excluding replicates with fewer than 20 observed MDCs (due to booststrap unreliability for small sample sizes), TODO/5 replicates with true neutral focal sites positive at $\alpha=0.05$.
Likewise, across the tested Sben/Gneu effect sizes, TODO/15 focal sites screened positive.
For true Sben/Gdel multilevel-selection effect conditions, we detected clade fitness statistics below the 50th percentile for all 5 $2\times$ effect size replicates.
However, sensitivity was insufficient to reliably detect smaller $1.1\times$ and $1.3\times$ effect sizes, with only TODO/10 replicates screened positive.

\subsection{Fitness Effect Screen: Dataset Size}

To assess the extent to which assay sensitivity could be improved by increases in dataset size, we performed additional replicates with a larger simulated population of 1.2 million agents.
Under these conditions, we were able to reliably detect true multilevel selection dynamics in TODO/TODO trials at the $1.3\times$ effect size treatment.
However, no multilevel selection effect was detected at the $1.1\times$ effect size with this larger population size.
Specificity remained fair for these larger population size trials, with false positive detections under Sben/Gneu conditions for only TODO/TODO trials.
Figure \ref{fig:samplesize:TODO} shows these results.

\subsection{Fitness Effect Screen: Environmental Fluctuations}

We next assessed robustness of comparator-based screens for negative MDC-associated fitness effects over populations under the influence of extrinsic, time-varying environmental factors.
For these experiments, we utilized Covasim configurations modeling vaccine rollout and public health policies over the course of the pre-omicron Covid-19 pandemic.
Due to significantly lower case counts within under this configuration, we conducted all trials at the larger 1.2 million host population size in order to achieve pathogen phylogeny dataset size comparable to baseline conditions.
Under these conditions, we found comparable sensitivity, with true positive detections in TODO/5 trials under $2\times$ effect size.
However, for smaller, $1.1\times$ and $1.3\times$ effect sizes, true Sben/Gdel dynamics were detected in only TODO/10 trials.
We also found specificity to be comparable to baseline conditions, with false positive detections in TODO/TODO trials.
Supplementary Figure \ref{fig:uk-summary-results} summarizes these results.

\subsection{Fitness Effect Screen: Adaptive Dynamics}

Finally, we assessed robustness of comparator-based screens for negative fitness effects to the presence of strong background inter-clade fitness differentials.
For these experiments, we extended our baseline Covasim configuration to incorporate introduction of alpha, beta, and gamma strains over the course of the pandemic.
In addition to partitioning MDC comparator sets by month, as was performed for other assays, for these trials we additionally assessed further partitioning by strain ID (e.g., wildtype, alpha, beta, etc.).
We include these strain-partitioned results in Supplementary Figure \ref{fig:multistrain-summary-results}, although did not find appreciable effects of this further clade-based partitioning.

As before, under month-by-month partitioning of comparator sets, we found reliable sensitivity under strong $2\times$ effect size.
Again, though, at smaller, $1.1\times$ and $1.3\times$ effect sizes, true Sben/Gdel dynamics were detected in only TODO/10 trials.
False positive detections occured in TODO/TODO trials.

\subsection{Mutation Frequency Screen}

\begin{figure}
\includegraphics[width=\linewidth]{binder/binder-2025-05-19-freqscreen-summary.ipynb/binder/teeplots/2025-05-19-freqscreen-summary/col=trt-name+hue=sig+kind=swarm+row=source-name+row_order=vanilla_uk-big_multistrain+viz=catplot+y=p-poisson+ext=.pdf}

\caption{%
\textbf{Mutation frequency screen results}.
\footnotesize
TODO.
}
\label{fig:mutfreq-screen}
\end{figure}


Figure \ref{fig:mutfreq-screen} shows focal mutation frequency detection outcomes across surveyed baseline, UK, and multistrain Covasim configurations.
Across all conditions, Sben treatments are reliably detected in all instances for $1.3\times$ and $2\times$ effect sizes.
In addition, no false-positive Sneu treatments are detected.
However, we do not find this test to be sensitive at the $1.1\times$ effect size level, with no Sben treatments detected.

Supplementary Figure \ref{fig:mutfreq-screen-counts} provides a raw view of focal MDC mutation counts observed across trials.
Notably, these counts help explain low sensitivity for small $1.1\times$ effect sizes, as actual MDC counts are highly overlapping with MDC counts under true neutral conditions.
